\cotefile{109}
\scanpage{10}
\begin{center}Reflexions Mathématique\end{center}

(The simplest thing \inserted{for me} is to write out a\deleted{\unsure{n}} \inserted{possible} alternative version as our photocopy of the original is not that clear now. I'll underline in red places where the words are changed from your original whether by Breen or by us.)

After a twelve years' silence, the time has come for the author's works of maturity, with his vision and style renewed. Here is \uline{a} day-by-day account of an explorer's \deleted{travel} \inserted{\uline{journey}} with occasional reflections on the travel \uline{itself} \unsure{((|(}as well as on the traveler and the \uline{manifold} world around \uline{him},) \inserted{thus} recapturing its genuine nature: \uline{that} of an impassioned adventure, rooted in life

In the two forthcoming volumes, under the common title ``Pursuing Stacks'' (the first volume of which is subtitled ``The Modelizing Story,''), the author sketches some themes of a vast sythesis of homotopical algebra, (co)homological algebra and topos theory -- ripe now for more than fifteen years \uline{but} never yet begun. The leading thread, tenacious and omnipresent (while often remaining implicit) comes from algebraic geometry, and from the intuitions (rich precise and provocat\uline{ive}ly fragmentary) \uline{surrounding} the notions of sheaves, stacks, $2$-stacks ``etc.'' on a topos.

Thes two volumes are the first in a planned, considerably wider series of ``Réflexions Mathématiques'' \uline{in which} the author intends to \uline{present}, among other things, some ``naive ideas'' about two ``new continents'', eager to be discovered and explored. One \uline{stems from}
% archive sheet number: 9
